\section{Capacity estimates and uniform compactness}

Insert a capacity test into the comparison potential. Modified comparison then bounds the capacity of one sublevel by the Monge--Amp\`ere mass of a slightly larger one. For an $L^p$ density, the capacity estimate of \cref{cor:lp-capacity-square} reverses this relation: a nonempty sublevel of fixed height has a definite capacity. The tail bound \eqref{eq:global-capacity-tail} bounds the depth of a normalized potential. Applied to differences of solutions, the same estimate upgrades $L^1$ compactness to uniform compactness. Write $\operatorname{Cap}=\operatorname{Cap}_\omega$ for the capacity of \cref{sec:hermitian-capacity}.

\Needspace{15\baselineskip}
\begin{lemma}\label{lem:thin-sublevel-capacity}
Let $u,v\in\PSH(X,\omega)\cap L^\infty(X)$, with $-H\leq v\leq0$ and $H\geq1$. Fix $0<\varepsilon<1$, and put
\[
 m=\inf_X(u-(1-\varepsilon)v),\qquad
 U_s=\{u<(1-\varepsilon)v+m+s\}.
\]
There are $\delta_\varepsilon>0$, depending only on $(X,\omega,\varepsilon)$, and $C_H>0$, depending only on $(n,H)$, such that
\begin{equation}\label{eq:thin-capacity}
 t^n\operatorname{Cap}(U_s)\leq C_H\int_{U_{s+t}}\omega_u^n
 \qquad(0<s,t,\ s+t<\delta_\varepsilon).
\end{equation}
Every $U_s$, $s>0$, has positive capacity.
\end{lemma}
\begin{proof}
Use the constant $B_0$ from \cref{thm:modified-comparison} and set
\[
 \delta_\varepsilon=\min\{1-\varepsilon,\varepsilon^3/(128B_0n^2)\}.
\]
For a capacity test $0\leq\rho\leq1$, put
\[
 \tau=\frac{t}{2(1-\varepsilon)(H+1)},\quad
 w=(1-\tau)v+\tau\rho,\quad
 m'=\inf_X(u-(1-\varepsilon)w).
\]
Then $0<\tau<1$ and
$0\leq(1-\varepsilon)(w-v)\leq t/2$, so $m-t/2\leq m'\leq m$. Consequently
\[
 U_s\subset V:=\{u<(1-\varepsilon)w+m'+s+t/2\}
 \subset U_{s+t}.
\]
Apply modified comparison to $u,w$ at height $s+t/2$. Its inverse factor is at most $(8/7)^n$. Moreover,
\[
 \omega_{(1-\varepsilon)w}
 \geq(1-\varepsilon)\tau\omega_\rho
 =\frac{t}{2(H+1)}\omega_\rho.
\]
Integration over $U_s\subset V$ and comparison over $V$ give \eqref{eq:thin-capacity}, with $C_H=[2(H+1)]^n(8/7)^n$, after taking the supremum over $\rho$. Each $U_s$ is nonempty by the definition of the infimum. \Cref{lem:strict-set-positive-volume} gives positive smooth volume, hence positive capacity by the constant test.
\end{proof}

\Needspace{11\baselineskip}
\begin{lemma}\label{lem:backwards-capacity}
In the preceding setting, suppose
\[
 \int_E\omega_u^n\leq A\operatorname{Cap}(E)^2\qquad(E\subset X\ \hbox{Borel}),
 \qquad A>0.
\]
For every $0<D<\delta_\varepsilon$,
\begin{equation}\label{eq:sublevel-capacity-lower}
 \operatorname{Cap}(U_D)\geq\frac{D^n}{4^nC_HA}.
\end{equation}
\end{lemma}
\begin{proof}
Write $a(s)=\operatorname{Cap}(U_s)^{1/n}>0$ and $K=(C_HA)^{1/n}$. The preceding lemma gives
\[
 t\,a(s)\leq K a(s+t)^2\qquad(s+t<\delta_\varepsilon).
\]
Suppose $4Ka(D)<D$. Starting with $s_0=D$, define
\[
 t_j=2Ka(s_j),\qquad s_{j+1}=s_j-t_j.
\]
We check positivity before using the sublevel inequality. If $a(s_i)\leq2^{-i}a(D)$ for $0\leq i\leq j$, then
\[
 \sum_{i=0}^j t_i\leq2Ka(D)\sum_{i=0}^j2^{-i}<4Ka(D)<D.
\]
Hence $s_{j+1}>0$, and the displayed sublevel inequality gives $a(s_{j+1})\leq a(s_j)/2$. Induction justifies every step, and $s_j\downarrow s_\infty\geq D-4Ka(D)>0$. Monotonicity gives
$0<a(s_\infty)\leq a(s_j)\leq2^{-j}a(D)$, a contradiction. Therefore $D\leq4Ka(D)$, which is \eqref{eq:sublevel-capacity-lower}.
\end{proof}

\begin{theorem}\label{thm:lp-potential-estimate}
There are constants $C_p,c_p>0$, depending only on $(X,\omega,\mathrm{d}V,p)$, such that, whenever
$u\in\PSH(X,\omega)\cap L^\infty(X)$, $\sup_Xu=0$, and
\[
 \omega_u^n=F\,\mathrm{d}V,\qquad F\in L^p(X,\mathrm{d}V),\quad F\geq0,
\]
we have
\begin{equation}\label{eq:lp-potential-bound}
 \|u\|_\infty\leq C_p(1+\|F\|_p),
 \qquad \|F\|_p\geq c_p.
\end{equation}
\end{theorem}
\begin{proof}
The zero density is excluded by \cref{cor:bounded-mass-positive}. Apply \cref{cor:lp-capacity-square,lem:backwards-capacity} with $v=0$, $H=1$, $\varepsilon=1/2$, and $A=C_p\|F\|_p>0$. Fix $D=\delta_{1/2}/2$. If $M=-\inf_Xu$, the resulting inequality is
\[
 \operatorname{Cap}(\{u<-M+D\})\geq c/\|F\|_p.
\]
If $M>D+1$, \eqref{eq:global-capacity-tail} bounds this capacity by $C/(M-D)$, giving the first assertion; the remaining case is immediate. Bounding the same capacity by $\operatorname{Cap}(X)$ gives the second assertion.
\end{proof}

If two uniformly bounded solutions remain separated at some point, the preceding estimate gives a fixed positive capacity to a sublevel of their difference. Small $L^1$ distance makes its volume small, while the uniform $L^p$ density bound makes its Monge--Amp\`ere mass small. This contradicts the sublevel estimate and upgrades $L^1$ compactness to uniform compactness. Put $p'=p/(p-1)$.

\begin{theorem}\label{thm:lp-uniform-compactness}
Suppose $u_j\in\PSH(X,\omega)\cap C(X)$, $\sup_Xu_j=0$, and
\[
 \omega_{u_j}^n=F_j\,\mathrm{d}V,\qquad F_j\geq0,\qquad \sup_j\|F_j\|_p<\infty
\]
for the fixed $p>1$. If $(u_j)$ is Cauchy in $L^1(X,\mathrm{d}V)$, then it is Cauchy uniformly. In particular, every such family has a subsequence converging uniformly to a continuous $\omega$-plurisubharmonic function.
\end{theorem}
\begin{proof}
By \cref{thm:lp-potential-estimate}, there is a common $H\geq1$ with $-H\leq u_j\leq0$. Suppose uniform Cauchy convergence fails. After exchanging indices when necessary, there are a number $a>0$ and arbitrarily large $j,k$ with
\[
 S:=\inf_X(u_k-u_j)\leq-a.
\]
\Needspace{10\baselineskip}
Choose $0<\varepsilon<1$ such that $\varepsilon H<a/4$, and choose a fixed $D>0$ with $2D<\delta_\varepsilon$ and $2D<a/4$. Apply the preceding sublevel lemmas with $u=u_k$, $v=u_j$. Since $v\leq0$,
$m=\inf(u-(1-\varepsilon)v)\leq S$. Thus
\[
 U_{2D}\subset\{u-v<S+2D+\varepsilon H\}
 \subset\{|u-v|>a/2\}.
\]
The capacity lower bound \eqref{eq:sublevel-capacity-lower}, with a common $A$ supplied by \cref{cor:lp-capacity-square}, gives $\operatorname{Cap}(U_D)\geq c>0$. On the other hand, \eqref{eq:thin-capacity}, H\"older and Chebyshev give
\[
 \begin{split}
 D^n\operatorname{Cap}(U_D)
 &\leq C_H\int_{U_{2D}}F_k\,\mathrm{d}V\\
 &\leq C_H\|F_k\|_p
       \left(\frac{2\|u_k-u_j\|_1}{a}\right)^{1/p'}
 \longrightarrow0.
 \end{split}
\]
This contradiction proves uniform Cauchy convergence. The normalized compactness theorem, \cref{thm:normalized-compactness}, supplies an $L^1$-Cauchy subsequence of any such family. Its uniform limit is continuous, and distributional convergence preserves $\omega$-plurisubharmonicity.
\end{proof}

\section{Construction of continuous solutions}

Smooth approximation gives solutions with unknown normalization scalars. A Gauduchon trace bounds those scalars above, and the lower bound for the density norm bounds them away from zero. The resulting uniform $L^p$ bound puts the potentials in the compact family just constructed, so their limit solves the prescribed equation.

\begin{proof}[of \cref{thm:kn-continuous}]
Choose smooth strictly positive $f_j\to f$ in $L^p(X,\mathrm{d}V)$. By \cref{thm:tw-smooth}, smooth existence gives
\[
 \omega_{u_j}^n=c_jf_j\,\mathrm{d}V,\qquad \sup_Xu_j=0,\qquad c_j>0.
\]
We first control the constants. For $n\geq2$, fix a Gauduchon metric $\gamma$. By \cref{prop:gauduchon-amgm},
\[
 c_j^{1/n}A_j\leq\int_X\omega\wedge\gamma^{n-1},
 \qquad A_j=\int_X(f_j\mathrm{d}V/\gamma^n)^{1/n}\gamma^n.
\]
Since $|x^{1/n}-y^{1/n}|\leq|x-y|^{1/n}$ for $x,y\geq0$, H\"older's inequality gives $A_j\to A>0$, where $A=\int_X(f\mathrm{d}V/\gamma^n)^{1/n}\gamma^n$. Thus $c_j$ is bounded above. In dimension one the same estimate is Stokes' identity and $A_j=\int_Xf_j\,\mathrm{d}V$.

The lower-density estimate in \cref{thm:lp-potential-estimate} gives
$c_j\|f_j\|_p\geq c_p$, so the constants are bounded away from zero. The densities $c_jf_j$ are now uniformly bounded in $L^p$, and the same theorem gives a common bound for $u_j$. Pass to a subsequence with $c_j\to c>0$ and $u_j$ Cauchy in $L^1$. By \cref{thm:lp-uniform-compactness}, $u_j\to u$ uniformly, with $u$ continuous and $\sup_Xu=0$.

Bounded Bedford--Taylor continuity gives $\omega_{u_j}^n\rightharpoonup\omega_u^n$, while $c_jf_j\to cf$ in $L^p$. Hence $\omega_u^n=cf\,\mathrm{d}V$, as required.
\end{proof}

For the exponential equation, the corresponding scalar is $\mathrm{e}^{\lambda\sup_Xu}$. Its upper bound comes from Jensen's inequality: weighting a Gauduchon probability by the root of the density relates the conserved trace to $\sup_Xu$. The lower-density estimate again gives the reverse bound.

\begin{lemma}\label{lem:exponential-gauduchon-mean}
Fix $\lambda>0$. Suppose smooth $u_j$ with $\omega_{u_j}>0$ solve
\[
 \omega_{u_j}^n=\mathrm{e}^{\lambda u_j}f_j\,\mathrm{d}V,
\]
where smooth strictly positive $f_j$ converge in $L^p$ to a nonnegative nonzero $f$. Then $M_j=\sup_Xu_j$ is bounded above and below, uniformly in $j$.
\end{lemma}
\begin{proof}
Absorb the fixed smooth positive ratio $\mathrm{d}V/\omega^n$ into both $f_j$ and $f$. This preserves their $L^p$ convergence, so we may use $\mathrm{d}V=\omega^n$ in the calculation. Choose $\rho>0$ smooth with $\ddc(\rho\omega^{n-1})=0$, using \cref{thm:gauduchon}; for $n=1$, take $\rho=1$. Put
\[
 \alpha=\int_X\rho\omega^n,\quad
 A_j=\int_X f_j^{1/n}\rho\omega^n,\quad
 \psi_j=u_j-M_j\leq0.
\]
Set $A=\int_X f^{1/n}\rho\omega^n$. The root inequality used in the preceding existence proof gives $A_j\to A>0$. Arithmetic--geometric mean and the conserved trace imply
\[
 \alpha=\int_X\omega_{u_j}\wedge\rho\omega^{n-1}
 \geq \mathrm{e}^{\lambda M_j/n}\int_X \mathrm{e}^{\lambda\psi_j/n}f_j^{1/n}\rho\omega^n.
\]
Apply Jensen's inequality to the probability measure
$A_j^{-1}f_j^{1/n}\rho\omega^n$. Taking logarithms yields
\begin{equation}\label{eq:gauduchon-exponential-mean}
 \lambda M_j\leq n\log(\alpha/A_j)
 +\frac{\lambda}{A_j}\int_X(-\psi_j)f_j^{1/n}\rho\omega^n.
\end{equation}
The normalized potentials $\psi_j$ have uniform $L^q$ bounds for every finite $q$, by \cref{lem:uniform-exponential}. With $q=np/(np-1)$, H\"older gives a uniform bound for the last integral, since
$\|f_j^{1/n}\|_{np}=\|f_j\|_p^{1/n}$. This proves the upper bound for $M_j$.

Finally, the density $F_j=\mathrm{e}^{\lambda u_j}f_j$ satisfies
$\|F_j\|_p\leq \mathrm{e}^{\lambda M_j}\|f_j\|_p$. The lower-density bound of \cref{thm:lp-potential-estimate}, applied to $\psi_j$, gives $c_p\leq \mathrm{e}^{\lambda M_j}\|f_j\|_p$, which bounds $M_j$ below.
\end{proof}

\Needspace{9\baselineskip}
\begin{proof}[of \cref{thm:nguyen-exp}]
Choose smooth strictly positive $f_j\to f$ in $L^p$. The smooth proper equation, \cref{thm:smooth-proper-existence}, gives smooth $u_j$ satisfying the required equations. The preceding lemma bounds $M_j$ and consequently bounds $F_j=\mathrm{e}^{\lambda u_j}f_j$ in $L^p$. The normalized potentials $\psi_j=u_j-M_j$ are uniformly bounded by \cref{thm:lp-potential-estimate}. Pass to a subsequence with $M_j\to M$ and $\psi_j$ Cauchy in $L^1$. \Cref{thm:lp-uniform-compactness} gives uniform convergence $u_j\to u\in\PSH(X,\omega)\cap C(X)$.

Uniform convergence and $f_j\to f$ in $L^p$ give
$\mathrm{e}^{\lambda u_j}f_j\to \mathrm{e}^{\lambda u}f$ in $L^p$. Bounded-product continuity therefore gives $\omega_u^n=\mathrm{e}^{\lambda u}f\,\mathrm{d}V$. Bounded exponential comparison, \cref{cor:bounded-exponential-comparison}, gives uniqueness, in fact among all bounded solutions.
\end{proof}
